\subsection{Reconstructing the harmonic completed triple and its finite detector}

The original proof of SECTIONS\_OVER\_Q A.134 and its mixed-exponent
input A.90 are absent from the recovered sources.  The following
construction replaces that missing dependency.  It uses the exact
coefficients and the local centres needed in the supplied A.137 and
A.139, but supplies a new, fully specified finite initialization.
It does not claim to reproduce the unavailable historical certificate.
Throughout this subsection put
\[
\begin{gathered}
 h=\tfrac12,\quad R=\tfrac{529}{382},\quad
 r=1119169,\quad Q=139r=155564491,\quad c=-\tfrac{338}{191},\\
 (a_0,a_1,a_h)=\left(\tfrac{529}{382},\tfrac{147}{382},
                     \tfrac{(13Q)^2}{191}\right),\qquad
 d=7\cdot17Q=18512174429,\quad \rho=\tfrac{34}{7Q},\\
 C^\circ:\quad T-c_i=a_iX_i^4\quad(i=0,1,h),\qquad
 (c_0,c_1,c_h)=(0,1,h),\quad X_0X_1X_h\ne0,\\
 S=\{2,3,7,13,17,23,139,191,r\}.
\end{gathered}
\tag{E.160.1}
\]

\begin{result}{E.160}{Proved}{A complete finite local and root-symbol initialization}
The curve $C^\circ$ has an ordinary point with $T\ne c$ over every
completion of $\mathbb Q$.  Its proper smooth model has genus $33$.
The finite centres may be chosen with
\[
 T_2=31/2,\quad T_{17}=1+a_1 17^4,\quad
 T_{139}=T_r=h+a_h,\quad
 T_p=R\quad(p=3,7,13,23,191).
\tag{E.160.2}
\]
There are choices of the fourth roots $X_{i,v}$ such that, for each
of the two displayed generators $D$ in its row,
\[
\begin{array}{c|c|c}
 i& D&\text{real sign of }X_i\\ \hline
 0&-1146,\ -382&-\\
 1&382,\ 191\cdot23\cdot13Q&+\\
 h&191,\ -3&+
\end{array}
\qquad
 \prod_v(X_{i,v},D)_v=1.
\tag{E.160.3}
\]
All $X_{i,p}$ are units except $v_{17}(X_1)=1$ and
$v_{457}(X_h)=1$.  Thus the rational normalizers
$(b_0,b_1,b_h)=(1,17,457)$ make $X_{i,p}/b_i$ a unit at every
finite prime and preserve all six product identities.
\end{result}
\begin{proof}
The three coordinate functions have independent divisors modulo
four.  Hence the geometric cover has degree $4^3=64$.  At each of
$0,1,h,\infty$ it has $16$ points of ramification index four.
Riemann--Hurwitz gives $2g-2=-128+4\cdot48=64$, so $g=33$,
with $64$ geometric boundary points.  At every prime outside $S$
it has tame good reduction.  The trace formula and Weil bound
\citep[(1.5.1), p.~275, and Theorem (1.6), p.~276]{RootDeligne1974}
give
\[
 \#C^\circ(\mathbb F_p)\ge p+1-66\sqrt p-64>0
 \qquad(p>5000).
\tag{E.160.4}
\]
Such an interior residue point lifts by Hensel's lemma, with all
three roots units.  Its $T$-coordinate can be chosen rational and
different from $c$: a sufficiently small rational perturbation
stays in the same three open fourth-power classes.

The retained exact certificate
\path{code/check_harmonic_audit.py} and its JSON output enumerate
all $661$ primes $p\le5000$ outside $S$.  For $660$ of them they
record the least residue $t$ for which all three nonzero ratios
$(t-c_i)/a_i$ are fourth powers.  At the remaining prime $457$
both $(h-c_0)/a_0$ and $(h-c_1)/a_1$ are fourth powers; take
$T_{457}=h+a_h457^4$.  Hensel gives the other two unit roots,
while $X_h=457$.  The same computation checks that both $191$
and $-3$ are squares modulo $457$.

At the odd primes in $S$, the exact rational centres in
(E.160.2) have all three ratio valuations divisible by four and
their unit residues fourth powers.  The sole nonunit root there
is $X_{1,17}=17$.  The certificate includes all these valuations,
unit residues and a fourth-root residue for each.  At $2$ the
three ratios are units with residues
\[
 17,\quad177,\quad209\pmod {256}.
\]
It records roots modulo $256$ in the classes $5,3,5$ modulo $8$,
respectively.  The derivative has valuation two, whereas the
error has valuation at least eight, so strong Hensel lifting
preserves those root classes.  At infinity use $T=R$ and the
signs in (E.160.3).

For completeness the symbol computation is specified without
depending on an opaque algebra package.  For odd $p$, with
$x=p^\alpha u$ and $D=p^\beta v$,
\[
 (x,D)_p=(-1)^{\alpha\beta(p-1)/2}
          (u/p)^\beta(v/p)^\alpha.
\]
At $2$ its exponent is
$((u-1)/2)((v-1)/2)+\alpha(v^2-1)/8+\beta(u^2-1)/8$
modulo two.  These are the tame-symbol formula and the explicit
dyadic formula in \citet[VIII, \S5.8--5.10, pp.~246--247,
and p.~248]{RootMilneCFT}.  Choose the odd-prime roots square
modulo the primes ramifying in their row's $D$, with one
exception: choose $X_{1,139}$ nonsquare.  The required choices
exist; at $r$, the residue $(h-1)/a_1$ is an eighth power, so
$X_{1,r}$ is square.  All these conditions are checked in the
certificate.  Its complete list of negative symbols is
\[
\begin{array}{c|c|c}
 i&D&\text{places with symbol }-1\\ \hline
0&-1146&2,\infty\\
0&-382&2,\infty\\
1&191\cdot23\cdot13Q&2,139.
\end{array}
\tag{E.160.5}
\]
All other symbols are positive.  Outside $S\cup\{457\}$ both
entries are odd-prime units.  At $457$ the only nonunit root
has both row characters split.  At $17$, both characters in
row $1$ are squares.  Finally division by the positive rational
$b_i$ preserves the product by Hilbert reciprocity
\citep[VIII, (5.10), p.~247]{RootMilneCFT}.
\end{proof}

\begin{result}{E.161}{Proved}{Exact root-stabilizer intersections for every refinement exponent}
Let $4\mid m$ and $a'_i=a_iw_i^4$, where $w_i\in\mathbb Q^\times$
is a unit at $3,13,23,191$.  For $j\ne i$ choose
\[
 \alpha_{ij}^m=(c_i-c_j)/a'_j,
\]
using the positive real root for a positive radicand and
$\zeta_{2m}$ times its absolute positive root otherwise.  Put
$K_i=\mathbb Q(\alpha_{ij}:j\ne i)$.
Fix an integer $M_0$ divisible by $8\prod_{p\in S\cup\{457\}}p$,
and put $F=\mathbb Q(\mu_{2mM_0})$, $L_i=FK_i$.  Then
\[
\begin{split}
K_0\cap F&=B_0=\mathbb Q(\sqrt{-1146},\sqrt{-382}),\\
K_1\cap F&=B_1=\mathbb Q(\sqrt{382},\sqrt{191\cdot23\cdot13Q}),\\
K_h\cap F&=B_h=\mathbb Q(\sqrt{191},\sqrt{-3}),
\end{split}
\qquad L_i^{\mathrm{ab}}=F.
\tag{E.161.1}
\]
Consequently, if $2mM_0\mid M$, every unit class modulo $M$
whose cyclotomic automorphism fixes $B_i$ contains arbitrarily
large primes at which the two radicands have simultaneous
$m$th roots.  No condition that this class be one modulo $m$
is required.
\end{result}
\begin{proof}
If a positive real number $\beta$ with $\beta^m\in\mathbb Q$
lies in an abelian extension of $\mathbb Q$, then
$\mathbb Q(\beta)$ is real and Galois.  All its conjugates are
real roots of $Z^m-\beta^m$, hence belong to $\{\beta,-\beta\}$.
Thus $\beta^2\in\mathbb Q$.

Apply this observation to absolute monomials in the two radicals.
The private valuations $v_{23}(a'_0)=2$,
$v_3(a'_1)=1$, $v_{13}(a'_h)=2$, and the common valuation
$v_{191}(a'_i)=-1$ determine exactly when a radical monomial
can belong to $F$.  In coordinates $(j=1,h)$ for row zero,
$(j=0,h)$ for row one, and $(j=0,1)$ for row $h$, its exponent
lattice is respectively
\[
\begin{split}
 R_0&=(m/2)\mathbb Z^2,\qquad R_h=(m/2)\mathbb Z^2,\\
 R_1&=\{(ma/4,mb/4):a,b\in\mathbb Z,\ a+b\equiv0\pmod2\}.
\end{split}
\tag{E.161.2}
\]
Indeed the valuation-one prime in rows $0,h$ forces its
exponent into $(m/2)\mathbb Z$; the valuation-two private prime
first forces $(m/4)\mathbb Z$, and the common prime $191$
then forces $(m/2)\mathbb Z$.  In row $1$, the two private
primes force multiples of $m/4$, and $191$ forces their
integer coefficients to have even sum.  Conversely these
conditions make every positive monomial a rational multiple
of the relevant fixed square roots.  The factors $w_j^4$
contribute rational numbers, since $4e_j/m$ is integral.

The actual, signed monomial values generating the relation
lattice give precisely the fields in (E.161.1).  In row $1$
the half powers of the two radicands have the same squareclass
$382$, and the quarter power of their product has squareclass
$191\cdot23\cdot13Q$; here
\[
 \frac{1}{a_0}\frac{1/2}{a_h}
       =\left(\frac{191}{23\cdot13Q}\right)^2.
\]
The signed half powers in rows $0,h$ give the other displayed
bases.  Their squareclasses are independent.

The perfect Kummer pairing
\citep[Theorem 5.30 and Remark 5.32, pp.~75--76]{RootMilneFT}
shows that monomials representing distinct cosets of $R_i$
are linearly independent over $F$.  Every rational polynomial
in the radicals is a sum of such monomials with coefficients
in the field generated by the relation values, namely $B_i$.
Its intersection with $F$ is therefore exactly $B_i$.
Moreover every intermediate field of the abelian extension
$L_i/F$ is spanned by the monomial eigenspaces whose characters
vanish on the corresponding subgroup.  If such a field is
abelian over $\mathbb Q$, the positive part of each generating
monomial has rational square by the first paragraph.  The
valuation calculation then puts that monomial in $F$.
This proves $L_i^{\mathrm{ab}}=F$.

The image of $\operatorname{Gal}(L_i/K_i)$ on $F$ has fixed
field $F\cap K_i=B_i$, hence is exactly
$\operatorname{Gal}(F/B_i)$.  An indicated residue class
therefore extends to an automorphism of $L_i$ fixing both
chosen radicals.  Also $L_i\cap\mathbb Q(\mu_M)=F$, since
the intersection is abelian over $\mathbb Q$.  The two
automorphisms agree on this intersection, and so define an
automorphism of the compositum.  Chebotarev
\citep[Theorem 8.31, p.~147]{CohMilneANT} gives arbitrarily
large primes in its conjugacy class.  Conjugation preserves
the cyclotomic residue and the existence of one simultaneously
fixed pair of roots.  Reduction at an unramified such prime
gives the asserted two residue roots.
\end{proof}

\begin{result}{E.162}{Proved}{A complete mixed-exponent construction with the exact harmonic initialization}
There are classes $A_0,A_1,A_h\in\widehat{\mathbb Q^\times}$
and ordinary $T_p\in\mathbb Q_p\setminus\{0,1,h,c\}$ at every
finite prime such that
\[
 (A_i)_p=\kappa_p(T_p-c_i),\qquad A_i\bmod4=[a_i].
\tag{E.162.1}
\]
The centres in (E.160.2) remain exactly those displayed.
All three global signs are positive and agree with the real
centre $R$.  Each coordinate has infinitely many private
primes of positive, ordinary integral valuation.  All finite
choices are effective; no infinite collection is claimed to
have been executed.
\end{result}
\begin{proof}
Start at exponent $n_0=4$ with coefficients
$a_{i,0}=a_i b_i^4$, where $(b_0,b_1,b_h)=(1,17,457)$.
Divide the roots in E.160 by $b_i$.  They are now units at
every finite prime and satisfy all six product identities.
The real signs are $(-,+,+)$.  Choose cofinal exponents
$n_j=4(j+1)!$.  We give one complete passage from $n=n_j$
to $m=n_{j+1}=n k$, with three fresh positive primes $q_i$.

The maintained coefficients are positive and differ from
the original $a_i$ by fourth powers of rationals that are
units at $3,13,23,191$.  The maintained curve has ordinary
local points with unit roots everywhere finite.  At the
first step its six symbol products are those of E.160;
after that step every individual symbol at a prime of $S$
will be positive, and the real roots will be positive.

Choose a finite set $\mathcal T$ containing $S\cup\{457\}$,
every prime dividing a coefficient or $mM_0$, all previously
frozen primes, and all primes through a bound $B_m$ for the
Weil estimate below.  Retain the existing centre $T_p$ at
each $p\in\mathcal T$ exactly.  Choose signs
$\varepsilon=(-,+,+)$ for the first step and $(+,+,+)$
thereafter.  At each such prime require
\[
 q_i\ \text{arbitrarily close to }X_{i,p}/\varepsilon_i.
\tag{E.162.2}
\]
All targets are units.  Choose the precision so that
$X_{i,p}/(\varepsilon_iq_i)$ has a $k$th root $Y_{i,p}$
arbitrarily close to one.  Openness of the local power map
near one, or strong Hensel lifting, gives this precision
also when $p\mid k$.

The congruences (E.162.2) are compatible with fixing $B_i$.
Here is the exact reciprocity check.  For any basis character
$D$ of $B_i$, the product of the local symbols of their unit
targets is
\[
 \prod_{p\in\mathcal T}(X_{i,p}/\varepsilon_i,D)_p
 =\frac{\prod_{p<\infty}(X_{i,p},D)_p}
        {\prod_{p<\infty}(\varepsilon_i,D)_p}=1.
\tag{E.162.3}
\]
Both products omitted outside $\mathcal T$ are trivial.
The numerator is the inverse real symbol by the maintained
product identity; the denominator is the same real symbol
by rational reciprocity, since $X_{i,\infty}$ has sign
$\varepsilon_i$.  To interpret this without assuming a prime
already exists, take any positive integer in the CRT class,
coprime to its modulus.  Hilbert reciprocity computes its
quadratic Dirichlet character as the displayed product.
Thus the cyclotomic automorphism of the class fixes both
generators of $B_i$.

Select $q_0,q_1,q_h$ in that order, each larger than $B_m$
and outside all old supports.  Before choosing $q_i$, use the
already updated coefficients in earlier coordinates and
the old coefficients in later coordinates.  E.161 gives a
prime satisfying (E.162.2) at which
\[
 (c_i-c_j)/a_j^{\rm current}\quad(j\ne i)
\]
have simultaneous $m$th roots.  At a previously chosen
$q_s$ in the same step, also impose
\[
 (\varepsilon_iq_i)^n\equiv
                 (c_s-c_i)/a_i^{\rm old}\pmod {q_s}.
\tag{E.162.4}
\]
This has a unit solution because, when $q_s$ was selected,
its right side had an $m$th root; a $k$th power of that root
is an $n$th root.  The added congruences do not alter
(E.162.3), since $q_s\notin S$ and both entries of every
row's Hilbert symbol are units there.  Enlarge the CRT
modulus by these primes and by $2mM_0$.  E.161 still applies.
Set
\[
 a_i^{\rm new}=a_i^{\rm old}(\varepsilon_iq_i)^n.
\tag{E.162.5}
\]
In particular the first negative factor is essential to
the chosen real-root compatibility in (E.162.3); its even
power keeps the coefficient positive.

At every $p\in\mathcal T$, the roots $Y_{i,p}$ just selected
satisfy $T_p-c_i=a_i^{\rm new}Y_{i,p}^{m}$, so the centre is
unchanged.  At a new prime $q_s$ put
$T_{q_s}=c_s+a_s^{\rm new}$ and $Y_s=1$.  In earlier
coordinates the required nonzero residue ratio is an
$m$th power by the root condition for $q_s$; in later
coordinates it is one by (E.162.4).  Hensel gives unit roots,
since $q_s\nmid m$.  These centres are positive and hence
different from $c<0$.

At any remaining prime the new curve has tame good reduction.
For the exponent $m$ its degree is $m^3$, its genus is
$g_m=m^3-2m^2+1$, and it has $4m^2$ geometric boundary points.
For example
\[
 B_m=(2g_m+4m^2+2)^2
\]
ensures $p+1-2g_m\sqrt p-4m^2>0$ when $p>B_m$.
Choose an interior residue point and lift it, with a rational
centre different from $c$.  All roots are units.  At the
primes in $S$ the already chosen roots are close enough to
one to be squares.  At all other primes the $D$ are units.
Thus all finite symbols for the new roots are positive;
choose all real roots positive at $T=R$.  This proves every
induction invariant.  Excluding all old support primes also
preserves the four valuations used in E.161.

At stage $j$ freeze, in addition, all primes at most $j$.
Every fixed finite prime therefore has its centre unchanged
from some stage onward.  The coefficients are a Cauchy
sequence in the full power completion: each ratio in
(E.162.5) is an $n_j$th power and $n_j$ is eventually divisible
by every integer.  Write its limit as $A_i$.  At a fixed
prime the maintained roots are units and
$a_{i,j}=(T_p-c_i)X_{i,p}^{-n_j}$.  Exponentiation by $n_j$
tends uniformly to one on $\mathbb Z_p^\times$, since every
finite quotient of that compact group has finite exponent.
Hence the local limit is the ordinary element $T_p-c_i$,
giving (E.162.1).  Fourth-power classes and signs were
preserved at every step.  A fresh $q_i$ acquires valuation
$n_j>0$ in just coordinate $i$, and is never used again.
This proves all support and convergence assertions.
\end{proof}

\begin{result}{E.163}{Proved}{The uniform proper Brauer obstruction and the missing completed coordinate}
The quaternion class
\[
 \alpha_C=(T-c,d)\in\operatorname{Br}(\overline C)[2]
\]
is defined on the proper smooth model and has invariant $1/2$
at every $2$-adic point and zero at every point over every
other completion.  In particular $\overline C(\mathbb Q)$
is empty, and no rational squareclass $b$ makes
$T-c=bY^2$, $Y\ne0$, soluble everywhere over $C^\circ$.
For the completed triple in E.162 there is no rational
squareclass, and hence no completed class, whose finite
localizations are the classes of $T_p-c$.
\end{result}
\begin{proof}
Put $a=a_0a_1a_h$ and $F(T)=T(T-1)(T-h)$.  The product map
$W=X_0X_1X_h$ takes $C^\circ$ to the affine curve
\[
 D^\circ:\quad W^4=F(T)/a,\qquad W\ne0,
 \quad T\ne0,1,h.
\tag{E.163.1}
\]
Its smooth proper model has genus three.  The following
identities are checked as exact rational equalities in the
certificate:
\[
 a=\frac{318039202706815117780707}{27871484},\qquad
 \frac{F(c)}a=-\rho^2,\qquad
 \frac{d}{2\rho}=\left(\frac{1088951437}{2}\right)^2.
\tag{E.163.2}
\]
At a point over $T=c$, its nonzero residue root $w$ gives
$\iota=w^2/\rho$, $\iota^2=-1$, and
$((1-\iota)w)^2=2\rho$.  Thus its residue field contains
$\sqrt d$.  The residue of $(T-c,d)$ vanishes there.
At infinity $T$ has pole order four on $\overline D$, so
that residue also vanishes; every other point has a unit
first entry.  The residue exact sequence for the Brauer
group of a smooth curve in characteristic zero
\citep[\S4.4, p.~20]{LitColliotThelene2015AWS}
therefore extends $(T-c,d)$ to $\alpha_D\in
\operatorname{Br}(\overline D)[2]$.  Its pullback is $\alpha_C$.

At $2$, the three coefficient valuations on $C$ are
$(-1,-1,0)$.  If $T$ is integral, the valuations of both
$T$ and $T-1$ would be $-1$ modulo four, impossible because
one of them is a unit.  If $v_2(T)<-1$, all three differences
have the same valuation, again incompatible with those
coefficient valuations modulo four.  Thus $v_2(T)=-1$.
Since $v_2(c)=1$, also $v_2(T-c)=-1$.  The unit $d$ is
$5$ modulo eight, so $\mathbb Q_2(\sqrt d)$ is the unramified
quadratic extension.  Its norm subgroup has even valuation,
giving invariant $1/2$.

At $7$, the coefficient valuations are $(0,2,0)$.
Negative valuation of $T$ is incompatible with the three
fourth-power equations.  For integral $T$, they force
$T\equiv1\pmod7$.  Since $c\equiv6\pmod7$, the first entry
is congruent to $2$, a square.  At $139$ and $r$, the
coefficient valuations are $(0,0,2)$; the same argument
forces $T\equiv h$.  The unit $h-c=867/382$ has squareclass
$1146$, a square at both primes by the certificate.
These are the three ramified odd primes of $d$ apart from $17$.

At $17$ all coefficients are units.  If $v_{17}(T)<0$,
that valuation is divisible by four and its leading unit
has the squareclass of $a_0$, which is square modulo $17$.
This makes the Hilbert symbol trivial.  If $T$ is integral,
the three ratios modulo $17$ must belong to
$\{0,1,4,13,16\}$.  Direct substitution, recorded by the
certificate, leaves only $T\equiv1$.  Then $T-c\equiv9$
is square.

At $3$, the valuations $(0,1,0)$ force $T$ integral and
$T\equiv1$, while $c\equiv2$; hence $T-c$ and $d$ are units.
At $13$, the valuations $(0,0,2)$ force $T\equiv h$,
whereas $c\equiv0$, again giving two units.  At $23$ and
$191$, $d$ itself is a local square, as certified.

At an odd prime outside $S$ it is useful to work on the
larger curve $D$, so that the argument will also apply to
its other local sheets below.  All displayed constants
are units.  If $v_p(T)<0$, the equation gives
$4v_p(W)=3v_p(T)$, so $4\mid v_p(T)=v_p(T-c)$.
If $T$ is integral and $T-c$ is a unit, the symbol is
trivial.  Otherwise reduction at $T=c$ gives a nonzero
root $w^4=-\rho^2$; the identities in (E.163.2) then make
$d$ square in the residue field and hence in $\mathbb Q_p$.
This proves vanishing on every ordinary local $D$-point
outside $S$.  At the real place $d>0$.

These calculations prove the assertion on the dense ordinary
locus.  Evaluation of a proper Brauer class is continuous,
and every point of a smooth curve over a local field has
ordinary points arbitrarily near it; thus the assertion
extends across all missing finite and boundary points.
The continuity and reciprocity facts are recalled in
\citet[\S3.2 and Theorem 3.4, pp.~14--15]{LitColliotThelene2015AWS}.
The invariant sum is $1/2$, excluding a rational point.
If $T-c=bY^2$ were everywhere soluble, evaluation would
instead be the local class $(b,d)$ everywhere.  Its real
value is zero regardless of $b$, and reciprocity contradicts
the same finite invariant sum.  The argument applies in
particular to the centres of E.162, which remain on the
original quartic curve because their coordinates have the
original fourth-power classes.  A putative completed new
coordinate would provide such a rational class modulo two.
\end{proof}

\begin{result}{E.164}{Proved}{The actual order-$4096$, class-eight harmonic detector}
For the completed triple of E.162 put
\[
 A=-4A_0A_1,\qquad B=4A_h^2,\qquad x_p=4T_p(1-T_p).
\tag{E.164.1}
\]
This is an ordinary completed tripod pair with both
coordinates of infinite support.  Write $\Pi$ for the full
geometric fundamental group of the tripod, $E_U$ for its
arithmetic fundamental group, and
\[
 H_D=\ker\bigl(\Pi\xrightarrow{\,2\kappa(t)+\kappa(1-t)\,}
                         \mathbb Z/8\bigr),\qquad
 K_D=\overline{H_D^2[H_D,H_D]}.
\]
Then $K_D$ is normal in $E_U$ and
\[
 \Pi/K_D\simeq(C_2)^8\rtimes C_{16},\qquad
 |\Pi/K_D|=4096,\qquad \operatorname{class}(\Pi/K_D)=8,
\tag{E.164.2}
\]
where $C_{16}$ rotates the eight coordinates.  There is no
section of $E_U/K_D\to G_{\mathbb Q}$ whose finite local
images are geometrically conjugate to those of the specified
$x_p$ and whose real image is ordinary or cuspidal.
In particular this excludes locally matching sections of
the full metabelian quotient and of the lower-central
quotient of class eight, and excludes a full Selmer section
with pair $(A,B)$.  It reconstructs A.139 for this exact
coefficient initialization and proves the finite exclusion
in A.137 as well.  The separate exact class-seventeen
calculation is reconstructed in E.165 below.
\end{result}
\begin{proof}
The elementary identities
\[
 1-4T(1-T)=4(T-h)^2,\qquad
 A^2B=64(A_0A_1A_h)^2\equiv64a^2\pmod{8}
\tag{E.164.3}
\]
give the asserted completed pair and its finite local
coordinates.  None of the $x_p$ is $0$ or $1$ because
$T_p\notin\{0,1,h\}$.  The private supports in coordinates
$0,1$ give infinite support of $A$, and those in coordinate
$h$ give twice their positive valuations in $B$.  The real
ordinary centre is $4R(1-R)=-77763/36481$.

The curve in (E.163.1) is, over the tripod, the finite
degree-eight cover
\[
 W^8=\frac{t^2(1-t)}{64a^2},\qquad
 t=4T(1-T),\qquad T=h-\frac{4aW^4}{t}.
\tag{E.164.4}
\]
Both compositions follow from $T(T-1)=-t/4$ and
$(T-h)^2=(1-t)/4$.  The exponent at $t=1$ is one, so the
cover is geometrically connected, with geometric subgroup
$H_D$.  Its degree-four model over the $T$-line has four
totally ramified boundary points, and genus three.

Take free meridians $x,y$ around $0,1$ and the Nielsen basis
$u=xy^{-2}$, $v=y$.  A Schreier basis of the index-eight
subgroup $H_D$ is
\[
 u_j=v^juv^{-j}\ (0\le j<8),\qquad z=v^8.
\]
Its quotient by $K_D$ is elementary abelian on these nine
independent generators.  Conjugation by $v$ rotates the
$u_j$ and fixes $z$; also $v^{16}=z^2=1$.
These relations give mutually inverse homomorphisms to
the semidirect product in (E.164.2): its independent
order-two classes are the eight lamps and $v^8$.
Thus the kernel is exactly $K_D$, not just a subgroup of it.
The Kummer character defining $H_D$ is multiplied by a
cyclotomic unit under arithmetic conjugation, so $H_D$
and its characteristic subgroup $K_D$ are arithmetic-normal.

Write the lamp module as
$L=\mathbb F_2[C_8]=\mathbb F_2[\tau]/(\tau^8)$,
$\tau=v-1$.  Its actual lower-central recurrence is
\[
 \Gamma_j(\Pi/K_D)=\tau^{j-1}L\quad(2\le j\le8),
 \qquad\Gamma_9=1.
\tag{E.164.5}
\]
The dimensions are $7,6,5,4,3,2,1$; the weight-eight
commutator is $\tau^7u_0=\sum_{j=0}^7u_j\ne0$.
Also $\Pi'\subset H_D$ implies $\Pi''\subset K_D$.
These prove the two asserted arithmetic quotient maps.

Suppose a section as in the statement exists.  The combined
mod-eight Kummer character descends to this quotient;
on $K_D$ its integral values are multiples of sixteen.
Its rational Kummer class has the same localization as
$64a^2$ at every finite prime by (E.164.3).  Such classes
are globally equal: valuations of their rational ratio
are divisible by eight, so that ratio is a sign times an
eighth power, and a negative sign is not an eighth power
in $\mathbb Q_3$.  Thus the pullback of the torsor in
(E.164.4) is trivial.  It fixes a geometric sheet, and,
after geometric conjugation, the section lands in the
actual subgroup $E_D/K_D$ of $E_U/K_D$.

Every local point on $D$ above $x_p$ has
$T=T_p$ or $T=1-T_p$.  Since $F(1-T)=-F(T)$, the second
sheet exists precisely when $-1$ is a fourth power in
$\mathbb Q_p$, that is, for odd $p\equiv1\pmod8$;
it does not occur at $2$.  Outside $S$, E.163 proves that
$\alpha_D$ vanishes at every local $D$-point.  Within $S$
the only possible second sheets are at $17$ and $r$.
For the fixed centres of E.160,
\[
\begin{array}{c|cc}
 p&T_p-c&1-T_p-c\\ \hline
17&9&8\\
r&867/382&867/382
\end{array}
\quad\text{in }\mathbb F_p.
\tag{E.164.6}
\]
All four entries are nonzero squares; at $r$ this is
$(1146/r)=1$.  Hence every local sheet above the specified
point has invariant $1/2$ at $2$ and zero elsewhere.

Here is the exact descent of the Brauer class to the finite
geometric quotient, without assuming an arithmetic splitting
of $E_D$.  Lift $\alpha_D$ by the Kummer sequence to
$H^2_{\mathrm{et}}(\overline D,\mu_2)$ and restrict to $D^\circ$.
There is a natural isomorphism
\[
 H^2(E_D,\mathbb F_2)\simeq
 H^2_{\mathrm{et}}(D^\circ,\mu_2).
\tag{E.164.7}
\]
Indeed the geometric affine curve has free profinite
fundamental group, its finite-coefficient cohomology
vanishes in degrees at least two, and its degree-one
cohomology is the character group.  The group and etale
Hochschild--Serre spectral sequences therefore have the
same geometric terms and comparison maps.  The required
geometric vanishing, spectral sequence and comparison are
\citet[Theorem 14.9, Remark 14.10 and Proposition 14.12,
pp.~96--97; Theorem 21.1, p.~130]{LitMilneLEC2013}.

Represent the corresponding class by a central extension
$1\to\mathbb F_2\to\widetilde E_D\to E_D\to1$.
It splits over the free geometric group $H_D$: lift a
free basis to obtain a continuous homomorphism
$j:H_D\to\widetilde E_D$.  For any $g\in\widetilde E_D$,
the maps $h\mapsto gj(h)g^{-1}$ and
$h\mapsto j(\overline g h\overline g^{-1})$ differ by a
character $H_D\to\mathbb F_2$.  Every such character kills
$K_D$.  Therefore $j(K_D)$ is normal and intersects the
central kernel trivially.  Quotienting gives
\[
\begin{gathered}
 1\longrightarrow\mathbb F_2\longrightarrow
 \widetilde E_D/j(K_D)\longrightarrow E_D/K_D\longrightarrow1,\\
 \widetilde E_D\simeq E_D\mathbin{\times}_{E_D/K_D}
                          \bigl(\widetilde E_D/j(K_D)\bigr).
\end{gathered}
\tag{E.164.8}
\]
For the second assertion, injectivity is the trivial
intersection just noted; surjectivity follows by correcting
any lift with $j(k)$, $k\in K_D$.  Compactness gives the
topological isomorphism.  Thus the class genuinely
inflates from $H^2(E_D/K_D,\mathbb F_2)$.

Pull this descended class back along the proposed global
section.  It is a rational Brauer class.  At a finite
prime the stipulated geometric conjugator can be lifted
to the full geometric group.  Since $K_D\subset H_D$,
the fixed finite sheet then gives a Galois-fixed sheet
of the actual cover over $x_p$, hence an actual point of
$D^\circ(\mathbb Q_p)$.  Naturality in (E.164.7)--(E.164.8)
identifies the local Brauer value with that point's
evaluation.  All its possible sheets were treated above.
At infinity an ordinary image has value zero because
$d>0$.  For a cuspidal image, its fixed inertia orbit
gives a real point on the proper normalization above
the cusp.  The proper cohomology class pulls back through
that point: the map to $\overline D$ kills cusp inertia.
Its value is again zero.  The global invariant sum is
therefore $1/2$, contradicting Brauer reciprocity.

For a full Selmer section with the completed pair,
injectivity of the ordinary local Kummer map fixes
the local point as $x_p$.  A cusp cannot have these
two complete local Kummer coordinates: at cusps
$0,1,\infty$ the regular-coordinate conditions respectively
force $B=1$, $A=1$, or $B/A=-1$ locally, whereas an ordinary
pair $(x_p,1-x_p)$ satisfies none of them.  Thus the finite
exclusion applies to it.  Finally the harmonic quartic
cover of A.137 dominates the cyclic product cover in
(E.164.4), so its geometric subgroup $H$ lies in $H_D$,
and $\overline{H^2[H,H]}\subset K_D$.  This gives an actual
map from its order-$2^{136}$ quotient to our detecting
quotient and proves its finite-section exclusion as well.
The independent claim about that larger quotient's exact
class seventeen is the subject of E.165.
\end{proof}

\begin{result}{E.165}{Proved}{The exact class seventeen of the large harmonic quotient}
Let $\Lambda=\ker(\Pi\to\mathbb Z/2)$ for the character
taking $x$ to zero and $y$ to one, and put
\[
 H=\ker(\Lambda\to\Lambda^{\mathrm{ab}}/4),\qquad
 K=\overline{H^2[H,H]},\qquad Q_C=\Pi/K.
\]
These are the actual geometric subgroups of A.137.
Then $|Q_C|=2^{136}$ and $Q_C$ has exact nilpotency class
seventeen.  In the convention $[a,b]=aba^{-1}b^{-1}$,
start with $[x,y]$ and commute successively with the fifteen
letters
\[
                  \texttt{xxxxyxyxyyyyyyy}.
\tag{E.165.1}
\]
The resulting weight-seventeen element is nonzero.
Consequently the arithmetic quotient of class seventeen
and the third-derived arithmetic quotient both map to
the detecting quotient of A.137.  Combined with E.160--E.164,
this independently reconstructs all the stated exclusions
of A.137, including its exact quotient class.
\end{result}
\begin{proof}
The index-two subgroup $\Lambda$ is free on
$x,yxy^{-1},y^2$.  Thus the deck group $G=\Pi/H$ is
\[
 G=(C_4)^2\rtimes C_8,
\]
where $x=(1,0,0)$, $y=(0,0,1)$ and odd elements of $C_8$
swap the two $C_4$ coordinates.  The multiplication is
literal: for $a,b,a',b'\in\mathbb Z/4$ and
$e,e'\in\mathbb Z/8$,
\[
 (a,b,e)(a',b',e')=
 \begin{cases}
 (a+a',b+b',e+e'),&e\text{ even},\\
 (a+b',b+a',e+e'),&e\text{ odd}.
 \end{cases}
\tag{E.165.2}
\]
The images of $x,yxy^{-1},y^2$ are three independent
order-four generators of the even subgroup.  Hence its
kernel on $\Lambda$ is exactly $H$, not merely a quotient
of its mod-four abelianization.  The group $G$ has order
$128$ and class three: its commutator subgroup is generated
by the difference of the two $C_4$ coordinates, conjugation
by $y$ negates that difference, and a second commutator
is its double, which is central.

Use the connected covering graph with vertex set $G$ and
directed edges $(g,x)$ from $g$ to $gx$ and $(g,y)$ from
$g$ to $gy$.  Let $R=\mathbb F_2[G]$.  Its chain complex is
\[
 C_1=R e_x\oplus R e_y\xrightarrow{\partial} C_0=R,
 \qquad
 \partial(g e_x)=gx-g,\quad\partial(g e_y)=gy-g.
\tag{E.165.3}
\]
The graph has $128$ vertices and $256$ edges.  A spanning
tree shows that $Z_1=\ker\partial$ has dimension $129$.
More explicitly, the loops from the identity along a tree
to a non-tree edge and back form a free Schreier basis of
$H$ and a cycle basis after reduction modulo two.
Consequently the path-chain map identifies
\[
              H/\overline{H^2[H,H]}\simeq Z_1.
\tag{E.165.4}
\]
This also proves the assertion for the full profinite
group: both sides of (E.165.4) are the finite mod-two
abelian quotients of that finite free Schreier basis.

The actual finite quotient is the image of the path map
\[
 \Pi\longrightarrow (R^2)\rtimes G,\qquad
 x\longmapsto(e_x,x),\quad y\longmapsto(e_y,y),
\tag{E.165.5}
\]
whose cocycle law is $P(uv)=P(u)+\overline uP(v)$.
Its kernel is contained in $H$ and is exactly $K$ there
by (E.165.4).  Thus the representation is faithful on
$Q_C$.  Its order is $128\cdot2^{129}=2^{136}$, and
conjugation on the kernel $Z_1$ is the left regular
$G$-action, independently of the path component.

Here is the full finite lower-central calculation.
Let $c_2=[x,y]$, and take the four weight-four seeds
\[
 [[c_2,x],x],\quad [[c_2,x],y],\quad
 [[c_2,y],x],\quad [[c_2,y],y].
\tag{E.165.6}
\]
Their vertex components are the identity since $G$ has
class three.  The fourth lower-central subgroup of the
free group is the closed normal closure of these seeds:
its third subgroup is normally generated by $[c_2,x]$
and $[c_2,y]$, and making both generators commute with
$x,y$ makes that whole third subgroup central.  Since
the kernel in (E.165.4) is elementary abelian, the image
$M_4=\Gamma_4 Q_C$ is therefore exactly the $R$-span of
their four path vectors.

For every $j\ge4$ the exact recurrence is
\[
             M_{j+1}=(x-1)M_j+(y-1)M_j.
\tag{E.165.7}
\]
Indeed $M_j$ lies in the abelian kernel, and its
commutators with $g\in G$ are $(1-g)M_j$.  The
augmentation ideal is the right ideal generated by
$x-1,y-1$; telescoping a word in $x,y$ proves this
also for inverse letters.  As $M_j$ is an $R$-submodule,
(E.165.7) is the entire commutator subgroup.

The retained exact integer-bit certificate
\path{code/check_harmonic_large_quotient.py} uses precisely
(E.165.2)--(E.165.7).  It enumerates $128$ vertices and
$256$ chain coordinates, not the elements of $Q_C$.
Gaussian elimination over $\mathbb F_2$ gives
\[
\begin{array}{c|rrrrrrrr}
 j&4&5&6&7&8&9&10&11\\ \hline
 \dim M_j&124&121&115&107&97&85&71&55
\end{array}
\]
\[
\begin{array}{c|rrrrrrr}
 j&12&13&14&15&16&17&18\\ \hline
 \dim M_j&41&29&19&11&5&2&0.
\end{array}
\tag{E.165.8}
\]
At every step it verifies that the computed space consists
of cycles and is invariant under $x$ and $y$.
It separately evaluates the individual commutator
(E.165.1), both by the kernel recurrence and by multiplying
its complete expanded word in (E.165.5).  Its vector is
\[
                \sum_{g\in G}g e_x\ne0.
\tag{E.165.9}
\]
The expanded unreduced word has length $196606$;
its SHA256 and the full dimension list are retained
in the JSON certificate.  Thus $\Gamma_{17}Q_C\ne1$
and $\Gamma_{18}Q_C=1$ are certified in the actual quotient.

Finally $\Lambda$ is stable under the actual arithmetic
action: its mod-two $y$-character is fixed, since the
cyclotomic multiplier is odd and conjugation vanishes in
abelianization.  Also $H$ is characteristic in $\Lambda$, and $K$ is
characteristic in $H$, so both are normal in the full
arithmetic group.  Also $\Pi'\subset\Lambda$,
$\Pi''\subset\Lambda'\subset H$, and
$\Pi'''\subset[H,H]\subset K$.  Together with
$\Gamma_{18}\Pi\subset K$, these give the asserted
actual arithmetic maps.  E.164 already proves the
section obstruction for $E_U/K$ by its map to $E_U/K_D$,
so all quotient consequences now follow.
\end{proof}
